\documentclass[11pt]{amsart}
% claims-sync: status=AUTHOR_REVIEW manifest=CLAIMS.md

\usepackage[T1]{fontenc}
\usepackage{lmodern}
\usepackage{amsmath,amssymb,amsthm,mathtools}
\usepackage{microtype}
\usepackage{booktabs}
\usepackage{enumitem}
\usepackage{xcolor}
\usepackage[hidelinks]{hyperref}
\usepackage[nameinlink,capitalize,noabbrev]{cleveref}

\setlength{\emergencystretch}{2em}

\newtheorem{theorem}{Theorem}[section]
\newtheorem{proposition}[theorem]{Proposition}
\newtheorem{lemma}[theorem]{Lemma}
\newtheorem{corollary}[theorem]{Corollary}
\newtheorem{question}[theorem]{Question}
\newtheorem{remark}[theorem]{Remark}
\newtheorem{definition}[theorem]{Definition}

\newcommand{\C}{\mathbb C}
\newcommand{\PP}{\mathbb P}
\newcommand{\Sym}{\operatorname{Sym}}
\newcommand{\End}{\operatorname{End}}
\newcommand{\tr}{\operatorname{tr}}
\newcommand{\Hess}{\operatorname{Hess}}
\newcommand{\Seg}{\operatorname{Seg}}
\newcommand{\Span}{\operatorname{Span}}
\newcommand{\boxt}{\mathbin{\boxtimes}}
\newcommand{\Nsp}{\mathcal N}
\newcommand{\Msp}{\mathcal M}

\title[Normalized Hessian Spaces and Kronecker Products]{Normalized Hessian Spaces and Kronecker Products\\of Symmetric Persistent Tensors}
\author[Adam Guetz]{Adam Guetz\textsuperscript{*}}
\thanks{\textsuperscript{*}Learned Response Labs. External-review preprint; the headline results are proof candidates pending human mathematical review. See ``A note on process'' at the end of this document.}
\date{August 18, 2026\\[0.35em]\normalfont\small
Ancillary code DOI: \href{https://doi.org/10.5281/zenodo.22004644}{\textcolor{magenta}{10.5281/zenodo.22004644}}}

\begin{document}

\begin{abstract}
We revisit the Kronecker-closure claim for symmetric persistent tensors in
Gharahi's recent preprint \emph{Kronecker Products of Symmetric Persistent
Tensors}.  The proof of its polarized Hessian identity establishes the desired
perfect-power determinant identity on tuples of decomposable polarization
vectors and then extends the proposed multilinear root by the universal
property of the tensor product.  This does not by itself extend the determinant
identity to arbitrary polarization tensors: the obstruction is precisely the
Segre ideal already noted earlier in that paper.

No counterexample to the closure theorem is produced here.  Instead we develop
a normalized polarized-Hessian space that gives several replacement results.
If one factor has a simultaneously strictly triangularizable normalized
Hessian space and the other has a nil normalized Hessian space, then their
Kronecker product is persistent.  If, in addition, the certificate of the first
factor factors through a single linear form, no nil-space hypothesis is needed
on the second factor.  A trace-zero Jordan-closure criterion implies the
required simultaneous triangularizability by Jacobson's generalized Engel
theorem for weakly closed sets.  These arguments give unconditional Kronecker
closure when one factor has dimension at most three, in every degree, and for
cubics when one factor has dimension at most four.

We also prove a Veronese-contraction lemma reducing the fixed-linear-form
condition in arbitrary degree, for each fixed dimension, to a quartic base
case.  Finally, exact symbolic computations on several hundred persistent
examples show Jordan closure with no observed failure, including tests on
Kronecker products formed in the full product tensor space.  The general
implication from persistence to the required normalized-Hessian structure
remains open.
\end{abstract}

\maketitle

\section{Introduction}

Persistent tensors were introduced as a class adapted to the substitution
method and tensor-rank lower bounds; in the symmetric setting they admit a
strong Hessian characterization \cite{GO25,GL24}.  Gharahi's recent preprint
\cite{G26} proposes an especially natural closure property: if
\[
  f\in \Sym^n V,\qquad g\in \Sym^n W
\]
are symmetric persistent tensors, then their Kronecker product
$f\boxt g\in\Sym^n(V\otimes W)$ is again persistent.

The purpose of this note is constructive.  We first isolate a logical issue in
the proof of the polarized Hessian identity in \cite[Proposition~10]{G26}.
The issue does not supply a counterexample, and exact symbolic tests continue
to support the stated closure theorem.  We then give a matrix-space mechanism
that proves the theorem under explicit structural hypotheses and yields several
unconditional subcases.

Let $r=n-2$.  By the Hessian characterization of Gharahi--Ottaviani, a
persistent form $f\in\Sym^n V$, $\dim V=d$, admits a nonzero multilinear
certificate
\[
 P_f:V^r\longrightarrow \C
\]
such that
\begin{equation}\label{eq:certificate-intro}
 \det H_f(A)=P_f(A)^d,\qquad A\in V^r,
\end{equation}
where $H_f(A)$ is the constant Hessian matrix of the quadratic partial
polarization of $f$ at the tuple $A$.
After choosing an admissible base tuple $B$ with $P_f(B)\neq0$, define
\[
\begin{aligned}
 M_f(A)&=H_f(B)^{-1}H_f(A),\qquad
 \lambda_f(A)=\frac{P_f(A)}{P_f(B)},\\
 N_f(A)&=M_f(A)-\lambda_f(A)I.
\end{aligned}
\]
The central object of this note is the normalized Hessian space
\[
 \Nsp_f(B)=\Span\{N_f(A):A\in V^r\}\subseteq\End(V).
\]
The principal sufficient conditions are:
\begin{itemize}[leftmargin=2em]
\item[(T)] $\Nsp_f(B)$ is simultaneously strictly triangularizable;
\item[(Nil)] every element of $\Nsp_f(B)$ is nilpotent;
\item[(TR)] every element of $\Nsp_f(B)$ has trace zero;
\item[(JC)] $XY+YX\in\Nsp_f(B)$ for all $X,Y\in\Nsp_f(B)$.
\end{itemize}
We make no general base-independence claim for (Nil), (TR), or (JC); the
sufficient conditions below are existential in the choice of an admissible
base.

The main conclusions are as follows.

\begin{enumerate}[label=(\roman*),leftmargin=2.3em]
\item If $f$ has (T) and $g$ has (Nil), then $f\boxt g$ is persistent
      (\cref{thm:TNil}).
\item If $f$ has (T) and its certificate factors through one fixed linear
      form in all polarization slots, then $f\boxt g$ is persistent for
      \emph{every} persistent $g$ (\cref{thm:factorized}).
\item (TR)+(JC) imply (T).  The only external input is Jacobson's generalized
      Engel theorem for weakly closed sets of nilpotent transformations
      (\cref{thm:jacobson}).
\item Consequently, Kronecker closure is unconditional when one factor has
      dimension $d\le3$, for arbitrary degree; for $n=3$ it is unconditional
      when one factor has $d\le4$ (\cref{cor:small-d,cor:cubic-d4}).
\item A fixed-linear-form certificate in all degrees reduces, for each fixed
      $d$, to a quartic rank-one problem (\cref{thm:veronese}).
\end{enumerate}

The general theorem of \cite{G26} is therefore not established by the arguments
in this note, but its unresolved content is considerably more structured than
the original Segre-to-global step suggests.

\section{The Segre obstruction in the polarized proof}

We use the conventions of \cite{G26}.  If
$f\in\Sym^n V$ and $g\in\Sym^n W$, then for decomposable polarization
vectors
\[
 U^{(i)}=v^{(i)}\otimes w^{(i)},\qquad 1\le i\le r=n-2,
\]
the differentiation identities of \cite[Proposition~7 and
Corollary~8]{G26} give
\begin{equation}\label{eq:decomp-H}
 H_{f\boxt g}(U^{(1)},\ldots,U^{(r)})
 =\frac1{n!}\,H_f(v^{(1)},\ldots,v^{(r)})
        \boxt H_g(w^{(1)},\ldots,w^{(r)}),
\end{equation}
where the right-hand side is now an ordinary matrix Kronecker product.
Therefore, for persistent $f$ and $g$,
\begin{equation}\label{eq:decomp-det}
 \det H_{f\boxt g}(U^{(1)},\ldots,U^{(r)})
 =\left[
   \frac1{n!}P_f(v^{(1)},\ldots,v^{(r)})
                P_g(w^{(1)},\ldots,w^{(r)})
  \right]^{d_1d_2}.
\end{equation}
The product inside brackets is bilinear in each pair
$(v^{(i)},w^{(i)})$, so the universal property of the tensor product does
indeed define a unique multilinear form
\begin{equation}\label{eq:Ph-def}
 P_{f\boxt g}:(V\otimes W)^r\longrightarrow\C
\end{equation}
whose value on decomposable tuples is the bracketed expression in
\eqref{eq:decomp-det}.

The issue is that the universal property extends the \emph{candidate root}
$P_{f\boxt g}$, not the determinant identity.  For each polarization slot,
$\det H_{f\boxt g}$ has degree $D=d_1d_2$ in the coordinates of
$U^{(i)}$, whereas $P_{f\boxt g}^D$ has the same degree.  Equality of these
polynomials on rank-one tensors does not determine them on $V\otimes W$.
Their difference may lie in the ideal of the Segre variety.

This is exactly the obstruction noted immediately before Proposition~10 in
\cite[Remark~9]{G26}: polynomial equality on
\[
 \Seg(\PP V\times\PP W)\subset\PP(V\otimes W)
\]
does not imply global polynomial equality because the Segre ideal, generated
by the $2\times2$ minors of the coordinate matrix, is nonzero.  The same issue
persists independently in each polarization slot.

Representation-theoretically, restriction to decomposable tensors detects the
Cartan component
\[
 \Sym^D V^*\otimes\Sym^D W^*
 \subset \Sym^D(V\otimes W)^*,
\]
while the other Cauchy summands vanish upon the corresponding Segre
restriction.  Thus a repair must explain why the determinant residual has no
component in this kernel; multilinearity of the proposed root alone does not do
so.

\begin{remark}
This observation does not imply that the polarized identity is false.  In fact,
the symbolic calculations described in \cref{sec:computations} repeatedly
confirm the identity at nondecomposable polarization tensors.  The point is only
that an additional structural argument is required.
\end{remark}

\section{Normalized polarized Hessian spaces}

Let $f\in\Sym^n V$ be persistent, $d=\dim V$, and $r=n-2$.  Write
\[
 H_f(A)=H_f(a_1,\ldots,a_r)
\]
for the Hessian matrix of the quadratic partial polarization
$\partial_{a_1}\cdots\partial_{a_r}f$.  Fix once and for all a certificate
$P_f$ satisfying \eqref{eq:certificate-intro}.

\begin{definition}[Admissible base and normalized Hessian space]
An \emph{admissible base tuple} is $B=(b_1,\ldots,b_r)\in V^r$ with
$P_f(B)\neq0$.  Set
\[
 G_f=H_f(B),\qquad
 M_f(A)=G_f^{-1}H_f(A),\qquad
 \lambda_f(A)=P_f(A)/P_f(B),
\]
and
\begin{equation}\label{eq:N-def}
 N_f(A)=M_f(A)-\lambda_f(A)I.
\end{equation}
The associated normalized Hessian space is
\[
 \Nsp_f(B)=\Span\{N_f(A):A\in V^r\}.
\]
\end{definition}

Both $H_f(A)$ and $P_f(A)$ are multilinear in the $r$ polarization slots, so
$N_f(A)$ is multilinear as well.  In particular, relative to a basis of $V$,
$\Nsp_f(B)$ is spanned by at most $d^r$ basis-tuple evaluations.

For later use we record the multilinear extension of
\eqref{eq:decomp-H}.

\begin{lemma}[Global matrix expansion]\label{lem:matrix-expansion}
Let $h=f\boxt g$.  For arbitrary
$U^{(i)}\in V\otimes W$, choose finite decompositions
\[
 U^{(i)}=\sum_{a_i}v^{(i)}_{a_i}\otimes w^{(i)}_{a_i}.
\]
For a multi-index $\alpha=(a_1,\ldots,a_r)$ put
\[
 A_\alpha=(v^{(1)}_{a_1},\ldots,v^{(r)}_{a_r}),\qquad
 C_\alpha=(w^{(1)}_{a_1},\ldots,w^{(r)}_{a_r}).
\]
Then
\begin{equation}\label{eq:global-matrix-expansion}
 H_h(U^{(1)},\ldots,U^{(r)})
 =\frac1{n!}\sum_\alpha H_f(A_\alpha)\boxt H_g(C_\alpha).
\end{equation}
\end{lemma}

\begin{proof}
Equation \eqref{eq:decomp-H} holds on a decomposable choice in each slot.
Both sides of \eqref{eq:global-matrix-expansion} are multilinear separately in
each $U^{(i)}$, so expanding each chosen decomposition gives the result.
Notice that no determinant is taken in this step; this is precisely why the
extension is legitimate.
\end{proof}

We now give the first replacement for the disputed Segre-to-global step.

\begin{theorem}[Triangularizable--nil criterion]\label{thm:TNil}
Let $f\in\Sym^n V$ and $g\in\Sym^n W$ be persistent.  Suppose there are
admissible base tuples $B$ and $C$ such that
\begin{enumerate}[label=(\alph*)]
\item $\Nsp_f(B)$ is simultaneously strictly triangularizable, and
\item every element of $\Nsp_g(C)$ is nilpotent.
\end{enumerate}
Then $f\boxt g$ is persistent.
\end{theorem}

\begin{proof}
Write $d_1=\dim V$, $d_2=\dim W$, $D=d_1d_2$, and choose arbitrary
polarization tensors $U^{(1)},\ldots,U^{(r)}\in V\otimes W$.  Using
\cref{lem:matrix-expansion} and factoring the base Hessians gives
\begin{equation}\label{eq:Xi-factor}
 H_{f\boxt g}(U)=\frac1{n!}(G_f\boxt G_g)\,\Xi(U),
\end{equation}
where
\[
 \Xi(U)=\sum_\alpha M_f(A_\alpha)\boxt M_g(C_\alpha).
\]
Choose a basis of $V$ in which every $N_f(A)$ is strictly upper triangular.
Since
$M_f(A)=\lambda_f(A)I+N_f(A)$, the matrix $\Xi(U)$ is block upper
triangular, with $d_1$ identical diagonal blocks
\begin{equation}\label{eq:Y-def}
 Y=\sum_\alpha \lambda_f(A_\alpha)M_g(C_\alpha).
\end{equation}
Consequently
\[
 \det\Xi(U)=\det(Y)^{d_1}.
\]
Write
\[
 \mu=\sum_\alpha\lambda_f(A_\alpha)\lambda_g(C_\alpha).
\]
Then
\[
 Y=\mu I+
 \sum_\alpha\lambda_f(A_\alpha)N_g(C_\alpha)
   =\mu I+\widetilde N,
\]
with $\widetilde N\in\Nsp_g(C)$.  By hypothesis $\widetilde N$ is nilpotent,
so $\det Y=\mu^{d_2}$ and therefore
\begin{equation}\label{eq:Xi-det}
 \det\Xi(U)=\mu^D.
\end{equation}

The multilinear root \eqref{eq:Ph-def} satisfies, after expansion of the same
decompositions,
\[
 P_{f\boxt g}(U)
 =\frac1{n!}\sum_\alpha P_f(A_\alpha)P_g(C_\alpha)
 =\frac{P_f(B)P_g(C)}{n!}\,\mu.
\]
Finally,
\[
 \det(G_f\boxt G_g)
 =\det(G_f)^{d_2}\det(G_g)^{d_1}
 =\big(P_f(B)P_g(C)\big)^D.
\]
Taking determinants in \eqref{eq:Xi-factor} and using
\eqref{eq:Xi-det} gives
\[
 \det H_{f\boxt g}(U)=P_{f\boxt g}(U)^D
\]
for arbitrary $U$.  The Hessian characterization of persistence
\cite[Theorem~2(c)]{GO25} completes the proof.
\end{proof}

A stronger one-sided conclusion is available when the certificate of the
triangularized factor has rank one in its polarization slots.

\begin{definition}[Fixed-linear-form certificate]
We say that $f$ satisfies \emph{(a-multi)} if
\begin{equation}\label{eq:a-multi}
 P_f(a_1,\ldots,a_r)=c\,\ell(a_1)\cdots\ell(a_r)
\end{equation}
for some $0\neq c\in\C$ and $0\neq\ell\in V^*$.
\end{definition}

When $P_f$ is symmetric, this is the polarized form of condition (a) in
\cite{GO25}, namely that the ordinary Hessian is a nonzero scalar times a
$d(n-2)$-th power of a linear form.

\begin{theorem}[Factorized-certificate criterion]\label{thm:factorized}
Let $f\in\Sym^nV$ and $g\in\Sym^nW$ be persistent.  Suppose $f$ satisfies
(a-multi) and, at an admissible base $B$, $\Nsp_f(B)$ is simultaneously
strictly triangularizable.  Then $f\boxt g$ is persistent, with no nil-space or
triangularizability hypothesis on $g$.
\end{theorem}

\begin{proof}
Use the notation of the proof of \cref{thm:TNil}.  Since (a-multi) holds and
$B=(b_1,\ldots,b_r)$ is admissible, define slotwise normalized forms
\[
 \widehat\ell_i(v)=\frac{\ell(v)}{\ell(b_i)}.
\]
For decompositions
$U^{(i)}=\sum_{a_i}v^{(i)}_{a_i}\otimes w^{(i)}_{a_i}$, set
\begin{equation}\label{eq:Wi}
 W^{(i)}=\sum_{a_i}\widehat\ell_i(v^{(i)}_{a_i})w^{(i)}_{a_i}\in W.
\end{equation}
The diagonal block \eqref{eq:Y-def} becomes, by multilinearity of $M_g$,
\[
 Y=M_g(W^{(1)},\ldots,W^{(r)}).
\]
Persistence of $g$ therefore gives
\[
 \det Y=\lambda_g(W^{(1)},\ldots,W^{(r)})^{d_2}.
\]
As in \cref{thm:TNil}, block triangularity yields
\[
 \det\Xi(U)=\lambda_g(W^{(1)},\ldots,W^{(r)})^D.
\]
On the other hand, the multilinear extension defining $P_{f\boxt g}$ gives
\[
 P_{f\boxt g}(U)
 =\frac{P_f(B)P_g(C)}{n!}
   \lambda_g(W^{(1)},\ldots,W^{(r)}).
\]
The same determinant calculation as before proves the required $D$-th power
identity globally.
\end{proof}

\section{Jordan closure and Jacobson's generalized Engel theorem}

The simultaneous triangularizability hypothesis in the preceding theorems is a
normal-form condition.  The following criterion replaces it with polynomial
membership identities in the normalized Hessian space.

\begin{proposition}[Trace plus Jordan closure implies nilpotence]
\label{prop:TRJCNil}
Let $\Nsp\subseteq\End(V)$ be a linear space over a field of characteristic
zero.  Assume
\[
 \tr X=0\quad\text{for all }X\in\Nsp,
 \qquad
 XY+YX\in\Nsp\quad\text{for all }X,Y\in\Nsp.
\]
Then every $X\in\Nsp$ is nilpotent.
\end{proposition}

\begin{proof}
Jordan closure is equivalent to closure under squaring, since
\[
 XY+YX=(X+Y)^2-X^2-Y^2.
\]
Hence $X^k\in\Nsp$ for every $k\ge1$: from $X^k\in\Nsp$ we obtain
$2X^{k+1}=XX^k+X^kX\in\Nsp$.  Therefore
$\tr(X^k)=0$ for all $k\ge1$.  Newton's identities imply that all non-leading
coefficients of the characteristic polynomial of $X$ vanish, so
$\chi_X(t)=t^{\dim V}$.  Thus $X$ is nilpotent.
\end{proof}

The remaining passage from pointwise nilpotence to simultaneous strict
triangularizability is false for arbitrary linear spaces of matrices.  Jordan
closure supplies exactly the extra multiplicative closure needed for a classical
theorem of Jacobson.

\begin{theorem}[Trace--Jordan criterion]\label{thm:jacobson}
Let $\Nsp\subseteq\End(V)$ be a finite-dimensional linear space over $\C$.
If $\Nsp$ is trace zero and Jordan closed, then $\Nsp$ is simultaneously
strictly triangularizable.
\end{theorem}

\begin{proof}
By \cref{prop:TRJCNil}, every element of $\Nsp$ is nilpotent.  Moreover
$\Nsp$ is a \emph{weakly closed} set in Jacobson's sense: for every
$X,Y\in\Nsp$,
\[
 XY+1\cdot YX\in\Nsp.
\]
Jacobson's generalized Engel theorem \cite{J52,J62} states that the associative
enveloping algebra $\mathcal A$ generated by a weakly closed set of nilpotent
linear transformations on a finite-dimensional vector space is nilpotent.
Thus $\mathcal A^m=0$ for some $m$.

For completeness, a nilpotent associative algebra of transformations is
simultaneously strictly triangularizable.  If $\mathcal A\neq0$, choose the
largest $k$ with $\mathcal A^kV\neq0$; then any nonzero
$v\in\mathcal A^kV$ satisfies $\mathcal A v=0$.  The induced algebra on
$V/\C v$ is again nilpotent.  Induction on $\dim V$ gives a complete invariant
flag on which every element of $\mathcal A$, hence every element of $\Nsp$,
acts strictly triangularly.
\end{proof}

Combining \cref{thm:TNil,thm:jacobson} gives a theorem-shaped strengthening of
the computational Jordan test.

\begin{corollary}\label{cor:TRJC-product}
Let $f$ and $g$ be persistent.  If, for some admissible bases, $\Nsp_f$ satisfies
(TR)+(JC) and $\Nsp_g$ satisfies (Nil), then $f\boxt g$ is persistent.  If
$f$ additionally satisfies (a-multi), the hypothesis on $g$ is unnecessary.
\end{corollary}

This isolates the following identity-shaped open question.

\begin{question}[Normalized Hessian Jordan closure]\label{q:JC}
Does every persistent symmetric tensor admit an admissible base for which its
normalized Hessian space is trace zero and Jordan closed?
\end{question}

A positive answer would imply the triangularizability mechanism needed above.
It is stronger than simultaneous triangularizability as a statement about
abstract matrix spaces, but it is also more algebraic: Jordan closure can be
expressed by explicit polynomial span-membership identities.

\section{Unconditional low-dimensional consequences}

The classifications in \cite[Theorem~5]{GO25} make the preceding criteria
unconditional in several nontrivial regimes.

\begin{corollary}[One factor of dimension at most three]\label{cor:small-d}
Let $f\in\Sym^n\C^{d_1}$ and $g\in\Sym^n\C^{d_2}$ be persistent.  If
$\min(d_1,d_2)\le3$, then $f\boxt g$ is persistent.
\end{corollary}

\begin{proof}
For $d=2$ and $d=3$, \cite[Theorem~5]{GO25} gives, up to $\mathrm{GL}$,
respectively the normal forms
\[
 x_0^{n-1}x_1,
 \qquad
 x_0^{n-1}x_2+x_0^{n-2}x_1^2.
\]
Their Hessian certificates satisfy (a-multi).  At the diagonal base
$(e_0,\ldots,e_0)$, a direct calculation gives
\[
 \Nsp_{x_0^{n-1}x_1}=\C E_{21}
\]
and, for the ternary normal form,
\[
 \Nsp\subseteq
 \Span\{E_{21}+c_nE_{32},\ E_{31}\}
\]
for a nonzero scalar $c_n$.  Both spaces are strictly lower triangular.
Thus the low-dimensional factor satisfies the hypotheses of
\cref{thm:factorized}; symmetry of the Kronecker product in its two factors
finishes the proof.
\end{proof}

For cubics there is a further dimension-four case.

\begin{corollary}[Cubics with one factor of dimension at most four]
\label{cor:cubic-d4}
Let $f\in\Sym^3\C^{d_1}$ and $g\in\Sym^3\C^{d_2}$ be persistent.  If
$\min(d_1,d_2)\le4$, then $f\boxt g$ is persistent.
\end{corollary}

\begin{proof}
For $n=3$, the certificate has one slot, so (a-multi) is automatic.  The cases
$d\le3$ follow from \cref{cor:small-d}.  For $d=4$, \cite[Theorem~5]{GO25}
puts a persistent cubic, up to $\mathrm{GL}$, in the form
\[
 f=a x_0^2x_3+b x_0x_1x_2+c x_1^3,\qquad ab\neq0.
\]
Its Hessian at $e_0$ is nonsingular, and with
$M(v)=H_f(e_0)^{-1}H_f(v)$ one computes
\[
 M(v)-v_0I=
 \begin{pmatrix}
 0&0&0&0\\
 v_1&0&0&0\\
 v_2&\frac{6c}{b}v_1&0&0\\
 v_3&\frac{b}{2a}v_2&\frac{b}{2a}v_1&0
 \end{pmatrix}.
\]
Hence the normalized Hessian space is simultaneously strictly lower triangular.
Apply \cref{thm:factorized}.
\end{proof}

These corollaries are independent of the disputed Segre-to-global inference in
\cite[Proposition~10]{G26}.

\section{A cubic algebra interpretation}

For cubics the Jordan-closure condition has an especially rigid interpretation.
Let $f\in\Sym^3V$ be persistent, choose $b\in V$ with $H_f(b)$ nonsingular, and
set
\[
 G=H_f(b),\qquad M(u)=G^{-1}H_f(u).
\]
Then $M(b)=I$.  Full symmetry of the cubic tensor implies
\begin{equation}\label{eq:comm-prod}
 M(u)v=M(v)u\qquad (u,v\in V).
\end{equation}
Indeed, pairing both sides with an arbitrary $w$ using the nondegenerate form
$G$ reduces \eqref{eq:comm-prod} to symmetry of the associated trilinear form.
Thus
\[
 u*v:=M(u)v
\]
defines a commutative bilinear product on $V$ with unit $b$.

\begin{proposition}\label{prop:cubic-assoc}
Assume the cubic normalized Hessian space is Jordan closed.  Then the product
$*$ above is associative.  Equivalently, $M(V)$ is a unital commutative
associative subalgebra of $\End(V)$ and $M(u)$ is left multiplication by $u$.
\end{proposition}

\begin{proof}
Jordan closure of $\Nsp_f$ is equivalent to Jordan closure of
$M(V)=\C I+\Nsp_f$.  Since persistence implies conciseness by definition \cite[Definition~1]{GO25}, the linear map
$M:V\to\End(V)$ is injective.  Hence for $u,v\in V$ there is a unique element
$q(u,v)$ with
\[
 M(q(u,v))=\frac12\big(M(u)M(v)+M(v)M(u)\big).
\]
Apply both sides to the unit $b$.  By \eqref{eq:comm-prod},
$q(u,v)=u*v$, so
\begin{equation}\label{eq:Jordan-L}
 M(u*v)=\frac12\big(M(u)M(v)+M(v)M(u)\big).
\end{equation}
Applying \eqref{eq:Jordan-L} to $w$ gives
\[
 2(u*v)*w=u*(v*w)+v*(u*w).
\]
Cyclically permuting $u,v,w$ yields three equations.  Subtracting any two shows
$(u*v)*w=u*(v*w)$ (over characteristic zero), hence associativity.
\end{proof}

This is slightly stronger than merely obtaining a Jordan algebra: in the Hessian
setting, Jordan closure of the regular multiplication operators forces the
underlying commutative product to be associative.  It suggests studying persistent
cubics through finite-dimensional local/Frobenius-type algebras rather than through
arbitrary nilpotent matrix spaces.

\section{A Veronese contraction reduction to quartics}

The factorized-certificate condition (a-multi) is known in the classified
low-dimensional cases but is open in general.  The following reduction localizes
its entire higher-degree content in degree four.

\begin{definition}
A form $f\in\Sym^m\C^d$, $m\ge3$, is \emph{certified} if there is a nonzero
multilinear form
\[
 P_f:(\C^d)^{m-2}\to\C
\]
such that $\det H_f=P_f^d$ identically in the polarization variables.
By \cite[Theorem~2(c)]{GO25}, certified and persistent are equivalent in the
present setting.
\end{definition}

\begin{theorem}[Veronese contraction lemma]\label{thm:veronese}
Fix $d$.  Suppose every certified quartic in $\Sym^4\C^d$ satisfies
(a-multi).  Then every certified form in $\Sym^m\C^d$, for every $m\ge3$,
satisfies (a-multi).
\end{theorem}

\begin{proof}
We induct on $m$.  The case $m=3$ is immediate because the certificate has one
slot; $m=4$ is the hypothesis.  Assume $m\ge5$, set $r=m-2\ge3$, and write
$P=P_f$.

Define the linear contraction map
\[
 \Phi:V\longrightarrow(V^*)^{\otimes(r-1)},\qquad
 \Phi(u)=\iota_uP=P(u,-,\ldots,-).
\]
Its kernel $K$ is a proper linear subspace because $P\neq0$.  For
$u\notin K$, commutativity of directional derivatives gives the exact matrix
identity
\[
 H_{\partial_u f}(v_2,\ldots,v_r)=H_f(u,v_2,\ldots,v_r).
\]
Hence $\partial_u f$ is certified and its certificate root is, up to a scalar
$d$-th root of unity, $\Phi(u)$.  By induction,
\begin{equation}\label{eq:Phi-veronese}
 \Phi(u)=c_u\,\ell_u^{\otimes(r-1)},\qquad u\notin K,
\end{equation}
for some $c_u\neq0$ and $\ell_u\neq0$.

We claim that all $\ell_u$ in \eqref{eq:Phi-veronese} are proportional.
Otherwise choose $u_1,u_2\notin K$ with independent $\ell_1,\ell_2$.
The affine line $u_1+t u_2$ meets the linear subspace $K$ in at most one point,
so choose $t\neq0$ with $u_t=u_1+t u_2\notin K$.  By linearity of $\Phi$,
\[
 c_t\ell_t^{\otimes p}
 =c_1\ell_1^{\otimes p}+t c_2\ell_2^{\otimes p},
 \qquad p=r-1\ge2.
\]
Restricting to the diagonal gives
\[
 a\ell_1(x)^p+b\ell_2(x)^p=e\ell(x)^p,
 \qquad abe\neq0.
\]
After absorbing $p$-th roots of $a,b$, the left side is $X^p+Y^p$ for
independent linear forms $X,Y$.  Over $\C$ this factors into $p\ge2$ distinct,
pairwise nonproportional linear factors, whereas the right side is a pure
$p$-th power.  Unique factorization gives a contradiction.

Thus the image of $\Phi$ is contained in the one-dimensional space
$\C\ell^{\otimes(r-1)}$ for a fixed $\ell$.  Since $V\setminus K$ spans $V$,
there is a linear form $c\in V^*$ with
\begin{equation}\label{eq:P-c-ell}
 P(u,v_2,\ldots,v_r)
 =c(u)\ell(v_2)\cdots\ell(v_r)
\end{equation}
for all arguments.

It remains to identify $c$ with $\ell$.  The determinant $\det H_f$ is symmetric
in the $r$ polarization slots, so for every $\sigma\in S_r$,
$(P\circ\sigma)^d=P^d$.  Since the polynomial ring is an integral domain,
$P\circ\sigma=\chi(\sigma)P$ for a character
$\chi:S_r\to\mu_d$.  Hence $P$ is either symmetric or alternating.  If it were
alternating, then $\iota_uP$ would be alternating in its $r-1\ge2$ remaining
slots, while \eqref{eq:P-c-ell} makes it symmetric; all contractions would
vanish, contradicting $P\neq0$.  Thus $P$ is symmetric.  Swapping the first two
slots in \eqref{eq:P-c-ell} gives
\[
 c(u)\ell(v)=c(v)\ell(u),
\]
so $c=\gamma\ell$ for some $\gamma\neq0$.  Therefore
$P=\gamma\ell^{\otimes r}$, which is (a-multi).
\end{proof}

\begin{remark}[The quartic base case]
The proof fails exactly at $r=2$: then $p=r-1=1$, and the degree-one Veronese
contains lines.  Moreover the alternating possibility is not eliminated by the
last contraction argument.  Thus the remaining base problem is genuinely:
for a certified quartic, prove that the bilinear certificate $P_f$ is symmetric
of rank one (and in particular not alternating).  This is not supplied by the
current low-dimensional classification for arbitrary $d\ge4$.
\end{remark}

\section{Exact symbolic evidence}\label{sec:computations}

The preceding results arose from exact symbolic experiments designed to test
possible repairs rather than merely to resample the Segre-restricted identity.
All arithmetic was over $\mathbb Q$ using SymPy; no floating-point tests were
used.  For each certified form, an admissible diagonal base was chosen when
available, the space $\Nsp_f$ was generated from all basis tuples, reduced to an
exact linear basis, and the following gates were checked:
\[
 \text{(TR)},\qquad
 \text{(JC)},\qquad
 \dim\bigcap_{N\in\Nsp_f}\ker N>0.
\]
Jordan closure was tested by exact rank membership for every unordered pair of
a reduced basis of $\Nsp_f$.

The corpus contained:
\begin{itemize}[leftmargin=2em]
\item ten named normal forms, including the binary $W_n$ family, the ternary
      normal forms, the classified persistent cubics in dimension four, and the
      quartic $x_0x_2^3+x_1x_2^2x_3+x_3^4$ from \cite[Example~28]{GO25};
\item all 156 certificate-passing forms found among the
      $\binom{35}{1}+\binom{35}{2}+\binom{35}{3}=7175$ coefficient-one sums of
      at most three distinct quartic monomials in four variables;
\item 151 accumulated certificate-passing quintics in four variables from the
      analogous partial search; and
\item four Kronecker products whose normalized Hessian spaces were constructed
      directly from the product tensor, including product dimensions
      $4,6,8,$ and $9$.
\end{itemize}
One named quartic is also present in the 156-form hunt, so these data comprise
321 gate executions but 320 distinct literal polynomial inputs before any
$\mathrm{GL}$-orbit reduction.  After literal deduplication, 2204 unordered
Jordan-pair membership tests remain.  No Jordan-closure failure was observed.

The negative controls were informative.  Six symmetric but nonpersistent forms
failed (JC) in every case.  The strictly triangular space
$\Span\{E_{12},E_{23}\}$ shows that simultaneous triangularizability does not
imply Jordan closure, since
\[
 E_{12}E_{23}+E_{23}E_{12}=E_{13}\notin
 \Span\{E_{12},E_{23}\}.
\]
Conversely, a certificate-satisfying bilinear toy family that is not the
polarized Hessian of a symmetric tensor passes (JC) but fails (TR); thus (JC)
alone is not sufficient.

A further empirical regularity is
\begin{equation}\label{eq:dim-observation}
 \dim\Nsp_f\in\{d-1,d\}
\end{equation}
for every distinct tested persistent input, even when $r\ge2$ and the a priori
spanning set has size $d^r$.  No proof of \eqref{eq:dim-observation} is known to
us.

There are important sampling limitations.  The quartic sweep uses coefficient
one only, the quintic block is partial, no $d=5$ hunt was completed, and the
known examples are heavily biased toward the fixed-linear-form condition.  The
computations should therefore be read as evidence for structural conjectures,
not as a substitute for proof.  Ancillary scripts and exact outputs will
accompany a public version of this note.

\section{Open problems and status of Kronecker closure}

The following problems now appear to isolate the main missing structure.

\begin{question}\label{q:persist-jc}
Does persistence imply (TR)+(JC) for a normalized polarized-Hessian space at
some admissible base?
\end{question}

By \cref{thm:jacobson}, a positive answer implies simultaneous strict
triangularizability and therefore supplies the matrix-space input to the
Kronecker repair.

\begin{question}[Quartic rank-one problem]\label{q:quartic}
If $f\in\Sym^4\C^d$ is certified and
\[
 \det H_f(u,v)=P_f(u,v)^d,
\]
must $P_f$ be a symmetric bilinear form of rank one?
\end{question}

A positive answer for each fixed $d$ implies, by \cref{thm:veronese}, the
fixed-linear-form condition in every degree in that dimension.

\begin{question}\label{q:dimN}
For a persistent symmetric tensor, is $\dim\Nsp_f\le d$?  More strongly, under
natural nondegeneracy hypotheses, must $\dim\Nsp_f\in\{d-1,d\}$?
\end{question}

There is also a direct route that bypasses normalized-Hessian triangularization.
Let
\[
 R_{f,g}(U)=\det H_{f\boxt g}(U)-P_{f\boxt g}(U)^D.
\]
The decomposable calculation shows that $R_{f,g}$ lies in the appropriate sum
of Segre ideals after restriction in each polarization slot.  A direct proof of
Kronecker closure could instead show that the persistence/Hessian relations on
$f$ and $g$ force every non-Cartan Cauchy component of $R_{f,g}$ to vanish.
This is an ideal-membership or syzygy problem rather than a density argument.

To summarize: the proof of \cite[Proposition~10]{G26} requires an additional
global argument, but the general closure theorem remains plausible.  The
results above establish it in substantial unconditional subcases and replace
the Segre-to-global step by explicit algebraic conditions whose general validity
can be attacked independently.

\section*{A note on process}
Learned Response Labs is a research collaboration between Adam Guetz and AI
systems (primarily ChatGPT and Claude).  This note originated in a reading of
Gharahi's preprint \cite{G26}; the gap analysis, the repair mathematics, the
drafting of this text, and the verification code were primarily generated by
AI systems (OpenAI GPT and Anthropic Claude) operating adversarially: results
produced by one system were independently rederived and adversarially audited by
another before internal acceptance, and every computation is exact rational
arithmetic reproducible from the ancillary bundle.  Adam designed and directed
the research program and made the circulation decision but did not perform
mathematical verification; no human has.  Accordingly, the results are
presented as \emph{proof candidates} that have undergone extensive
computational and cross-model audit but no independent human mathematical
verification.  We invite external review and particularly welcome corrections
to any argument or computation.

\section*{Acknowledgments}
The starting point for this note is Masoud Gharahi's preprint \cite{G26} and the
Hessian characterization developed with Giorgio Ottaviani in \cite{GO25}.  I am
grateful for the clarity of those formulations, which makes the obstruction and
the matrix-space alternatives unusually explicit.  This is a preliminary note;
comments and corrections are welcome.

\begin{thebibliography}{99}

\bibitem{G26}
M.~Gharahi,
\emph{Kronecker Products of Symmetric Persistent Tensors},
arXiv:2608.11182v1, 2026.

\bibitem{GO25}
M.~Gharahi and G.~Ottaviani,
\emph{Symmetric Persistent Tensors and their Hessian},
arXiv:2510.07404v1, 2025.

\bibitem{GL24}
M.~Gharahi and V.~Lysikov,
\emph{Persistent Tensors and Multiqudit Entanglement Transformation},
Quantum \textbf{8} (2024), 1238.

\bibitem{J52}
N.~Jacobson,
\emph{Une g\'en\'eralisation du th\'eor\`eme d'Engel},
C. R. Acad. Sci. Paris \textbf{234} (1952), 679--681.

\bibitem{J62}
N.~Jacobson,
\emph{Lie Algebras},
Interscience Tracts in Pure and Applied Mathematics, No.~10,
Interscience Publishers, New York--London, 1962.

\end{thebibliography}

\end{document}
